\section{Hai con trỏ: tổng 2 phần tử bằng $K$ [Trung bình]}
\subsection{Phân tích \& Thiết kế (M0)}
\begin{itemize}
    \item \textbf{Input:}
    File \texttt{input.txt} chứa:
    \begin{itemize}
        \item Số nguyên $n$
        \item $n$ số nguyên đã sắp xếp tăng
        \item Số nguyên $K$
    \end{itemize}

    \item \textbf{Xử lý:}
    \begin{itemize}
        \item \textbf{Case biên:}
        \begin{itemize}
            \item $n = 0 \rightarrow$ in \texttt{NONE}
            \item $n = 1 \rightarrow$ in \texttt{NONE}
            \item Chỉ in cặp khi $i < j$
            \item Nếu $a[i] = a[j] = K/2$ thì vẫn hợp lệ khi $i \ne j$
        \end{itemize}

        \item \textbf{Module 1 (Đọc file):}
        \begin{itemize}
            \item Hàm \texttt{readInput(path, vector<long long>\&, long long\&)}
            \item Kiểm tra mở file (\texttt{fail} $\rightarrow$ \texttt{cerr + return false})
        \end{itemize}

        \item \textbf{Module 2 (Tìm các cặp):}
        \begin{itemize}
            \item Dùng kỹ thuật \texttt{Two Pointers}
            \item \texttt{left = 0, right = n-1}
            \item Nếu tổng $= K$ $\rightarrow$ ghi cặp, tăng trái giảm phải
            \item Nếu tổng $< K$ $\rightarrow$ tăng \texttt{left}
            \item Nếu tổng $> K$ $\rightarrow$ giảm \texttt{right}
        \end{itemize}

        \item \textbf{Module 3 (Ghi file):}
        \begin{itemize}
            \item Mỗi cặp ghi:
\begin{lstlisting}
i j a[i] a[j]
\end{lstlisting}
            \item Nếu không có cặp nào ghi \texttt{NONE}
        \end{itemize}
    \end{itemize}

    \item \textbf{Output:}
\begin{lstlisting}
i j a[i] a[j]
\end{lstlisting}
Hoặc:
\begin{lstlisting}
NONE
\end{lstlisting}
\end{itemize}

\subsection{Mã nguồn C++11 (Tách module)}
\begin{lstlisting}
#include <iostream>
#include <vector>
#include <fstream>

using namespace std;

struct PairAns {
    int i, j;
    long long x, y;
};

bool readInput(const string& path, vector<long long>& a, long long& K) {
    ifstream fin(path);
    if (!fin) return false;

    int n;
    if (fin >> n) {
        a.resize(n);
        for (int i = 0; i < n; ++i) fin >> a[i];
        fin >> K;
    }
    return true;
}

vector<PairAns> findPairs(const vector<long long>& a, long long K) {
    vector<PairAns> ans;
    int n = (int)a.size();

    int left = 0, right = n - 1;

    while (left < right) {
        long long sum = a[left] + a[right];

        if (sum == K) {
            ans.push_back({left, right, a[left], a[right]});
            ++left;
            --right;
        }
        else if (sum < K) {
            ++left;
        }
        else {
            --right;
        }
    }

    return ans;
}

bool writeOutput(const string& path, const vector<PairAns>& ans) {
    ofstream fout(path);
    if (!fout) return false;

    if (ans.empty()) {
        fout << "NONE";
    } else {
        for (size_t i = 0; i < ans.size(); ++i) {
            fout << ans[i].i << " "
                 << ans[i].j << " "
                 << ans[i].x << " "
                 << ans[i].y << "\n";
        }
    }
    return true;
}

int main() {
    vector<long long> a;
    long long K;

    if (!readInput("input.txt", a, K)) {
        cerr << "Loi mo file input!\n";
        return 1;
    }

    vector<PairAns> ans = findPairs(a, K);

    if (!writeOutput("output.txt", ans)) {
        cerr << "Loi mo file output!\n";
        return 1;
    }

    return 0;
}
\end{lstlisting}

\subsection{Kế hoạch kiểm thử (Test Cases)}

\textbf{Test Case 1 (Thông thường)}

\textbf{Input (input.txt):}
\begin{lstlisting}
6
1 2 3 4 5 6
7
\end{lstlisting}

\textbf{Expected Output:}
\begin{lstlisting}
0 5 1 6
1 4 2 5
2 3 3 4
\end{lstlisting}

\vspace{0.5cm}

\textbf{Test Case 2 (Không có cặp)}

\textbf{Input (input.txt):}
\begin{lstlisting}
4
1 2 3 4
20
\end{lstlisting}

\textbf{Expected Output:}
\begin{lstlisting}
NONE
\end{lstlisting}

\vspace{0.5cm}

\textbf{Test Case 3 (Case biên -- giá trị giống nhau)}

\textbf{Input (input.txt):}
\begin{lstlisting}
4
5 5 5 5
10
\end{lstlisting}

\textbf{Expected Output:}
\begin{lstlisting}
0 3 5 5
1 2 5 5
\end{lstlisting}